% Generated by report/make_restated.py from theory/A1_lower_bound.tex: do not edit (edit the note and rerun).
\begin{restated}{Remark}{rem:rho}{A1}{0}{sharpness of the slow-mode bound, and the exhausted chain}
The constants in \eqref{ineq:rho} are not sharp: the proof loses a factor $m/S(\theta_k)$ in $S(\theta)\le m$ (for every mode, the first one included, $S(\theta_k)=\frac m2-\frac{\sin(m\theta_k)\cos((m+1)\theta_k)}{2\sin\theta_k}\approx\frac m2$, exactly $\frac{m+1}2$ for $\zeta=2$ and $\frac{2m+1}4$ for $\zeta=1$; for $m=200$, $S(\theta_k)/m\in[0.50,0.61]$ along the instance, the largest values being those of the first mode), a factor up to $4$ in $(\lambda_k+s)^2\le4s^2$, and the rounding in $K_0$. The exact $\rho$ is computable from the $m\times m$ tridiagonal eigenproblem and is the quantity to report in the experiments; for $n=100$, $m=200$ and $\kappa\in[10^3,3\cdot10^5]$ (so that $y\in[2.3,40.5]$; $4001$ log-spaced values of $\kappa$, completed by the jump points of both quantities) the exact $\rho$ ranges from $0.13$ to $0.35$, against $0.013$--$0.046$ for the bound \eqref{ineq:rho}, and the ratio of the bound to the exact $\rho$ lies between $0.06$ and $0.21$. The exact $\rho$ is not monotone in $\kappa$: it drops each time an eigenvalue $\lambda_k$ crosses $s$ (one mode leaves the slow set); its largest values occur near $\kappa=10^3$ and its smallest ones ($\approx0.13$) just above $\kappa\approx2.2\cdot10^5$, where the third mode leaves the slow set. For the infinite chain ($m=\infty$) the spectral measure of $e_1$ for $A$ is $d\nu(\lambda)=\frac1{2\pi}\sqrt{\lambda(4-\lambda)}\,d\lambda$ on $[0,4]$ (the orthogonal polynomials of the recursion of Step~1 are Chebyshev polynomials of the second kind), $\norm{x^\ast}^2=\int_0^4(\lambda+s)^{-2}d\nu(\lambda)$, and the slow fraction $\int_0^s(\lambda+s)^{-2}d\nu/\int_0^4(\lambda+s)^{-2}d\nu$ tends, as $s\to0$, to $\frac{\int_0^1\sqrt u\,(1+u)^{-2}du}{\int_0^\infty\sqrt u\,(1+u)^{-2}du}=\frac{\pi/4-1/2}{\pi/2}=\frac12-\frac1\pi\approx0.18$; this is the value approached by the exact $\rho$ when $m\sqrt s\to\infty$ and then $s\to0$, and it is not used in the proofs. At the other extreme, when $\kappa\to\infty$ with $m$ fixed ($s\to0$, $y\to0$), eventually no eigenvalue of $A_m(\zeta_Q)$ lies below $s$ and the slow set is empty; the relevant quantity is then the single-mode overlap $\delta:=\inner{x^\ast}{v_1}^2/\norm{x^\ast}^2$ of Corollary~\ref{cor:saga}(a) with $\nu'=\mu_F$ (the eigenspace of $\mu_F$ is spanned by $E_1v_1,\dots,E_nv_1$, so $\norm{P_{\mu_F}(\theta_0-\etastar)}^2=n\inner{x^\ast}{v_1}^2=\delta\norm{\theta_0-\etastar}^2$). Since $\zeta_Q\to1$, $x^\ast=(A_m(\zeta_Q)+sI)^{-1}e_1\to A_m(1)^{-1}e_1$ and $v_1(\zeta_Q)\to v_1(1)$ (simple eigenvalue, eigenvectors normalized with positive first coordinate), so $\delta\to\delta_\infty(m):=\inner{A_m(1)^{-1}e_1}{v_1(1)}^2/\norm{A_m(1)^{-1}e_1}^2$. In fact $A_m(1)\mathbf 1=e_1$ (row sums), so $A_m(1)^{-1}e_1=\mathbf 1$, and $v_1(1)$ is the normalized vector $(\sin\frac{j\pi}{2m+1})_{j\le m}$ (Lemma~\ref{lem:chain}), whose entries are positive; hence $\delta_\infty(m)=\inner{\mathbf 1}{v_1(1)}^2/m=\cot^2\big(\frac{\pi}{2(2m+1)}\big)/(m(2m+1))>0$, which is $0.81$ for $m=200$ and tends to $8/\pi^2$. The overlap $\delta$ approaches this value only once $s\ll\lambda_1$: for $m=200$ and $n=100$, $\delta\approx0.18$ at $\kappa=10^6$ (one mode is still slow there, so $\delta=\rho$), $0.56$ at $10^7$ and $0.74$ at $10^8$.
\end{restated}
